% ======================= 摘要页（第 1 页） =======================
\begin{center}
{\zihao{3}\heiti 基于共表达图与风格条件流的啤酒风味预测、归因与生成}
\end{center}

\vspace{2pt}
\begin{center}
{\zihao{4}\heiti 摘\hspace{1em}要}
\end{center}

\vspace{-2pt}
啤酒风味由两百余种化合物在极宽浓度域内的协同与拮抗共同决定，单一化合物的气味活度
值无法解释混合效应。本文以 250 款商业啤酒的 231 个化合物指标与 50 维 0--6 分制感官
评分为对象，按\textbf{数据审计、感官预测、化合物归因、风味生成}四步推进，每一问都
给出所用模型、具体做法与量化结果，全部数字可由随附源码一次复现。

\textbf{针对问题一}，本文采用\textbf{多任务图扩散网络 MT-GNN 与非负岭堆叠集成}。
做法上，先用校准浓度表反证出对数变换底数为 10（斜率 1.0000、$R^2=1.0000$），并识别
出 231 列中有 \textbf{33 列实为原始浓度}、\textbf{102 列存在检出限左删失}、73 个缺失
单元，据此构造缺失指示与封底指示两类变量；再把化合物相关矩阵转为\textbf{共表达图
扩散算子} $\bm{S}=(1-\alpha)\bm{I}+\alpha\tilde{\bm{A}}$ 与原始特征拼接，送入双分支
共享隐层网络一次输出 50 维，最后用非负岭回归对十个基模型做堆叠。结果表明，在按风格
分层的 5 折 $\times$ 3 重复交叉验证下，\textbf{集成模型平均 $R^2=0.268$}，高于文献
梯度提升基线的 0.21 与 0.22，\textbf{50 维中有 25 维超过该基线}，最高单维
$R^2=0.74$（醇感），最难的总体喜好维度 $R^2=0.125$，为文献同项的 1.4 倍；
\textbf{MT-GNN 是十个模型中唯一 50 维全为正者}（最低 $R^2=0.0005$），说明共享隐层
把统计强度传导到了弱信号维度。消融实验如实表明，缺失与封底指示、图扩散的边际增益
均在噪声范围内（$|\Delta|<0.01$），\textbf{增益主要来自多任务结构与集成而非单一
特征工程}。最终交付 $250\times50$ 的 0--6 分制评分表、整数分级表，以及实测覆盖率
为 7.5\% 的预测区间。

\textbf{针对问题二}，本文采用\textbf{贪心模块度优化与五路证据秩融合}。做法上，先在
$|\rho|\ge0.45$ 的共表达图上做社区发现，并剔除会挤占名额的理化参数代理变量，得到
34 个社区、\textbf{模块度 $Q=0.549$}，其中 6 个主社区覆盖 200 个化合物且与酿造生化
通路一一对应；再对每个感官维度并行计算\textbf{秩相关、SHAP、置换重要性、ElasticNet
稳定性选择、互信息}五路排序，加权融合后取前三名。稳定性用\textbf{bootstrap 入选
频率}与\textbf{跨方法 Jaccard 系数}双指标度量，化学合理性用阈值—OAV、化学类别、
代谢来源、社区一致性四条判据核验，并以 \texttt{6.csv} 交叉相关矩阵作独立外部验证。
结果表明，\textbf{Top-1 的 bootstrap 入选频率均值达 0.72}（随机基线 0.33），
\textbf{76\% 的维度其 Top-3 至少两项落在同一社区}，19/68 个可算 OAV 的条目
满足 $\mathrm{OAV}>1$，据此识别出里那醇、香叶醇、甲基庚烯酮、乙酸乙酯、DMS 等
公认关键致香物，统计证据与化学语义相互印证。

\textbf{针对问题三}，本文采用\textbf{风格条件潜因子流 SCLF 与带约束的高斯最优传输}。
做法上，先用并行分析定阶取 $K=14$ 个潜因子，以 6 层仿射耦合流建模风格条件密度
$p(\bm{z}\mid s)$，采样时按该风格缺陷 OAV 的 95\% 分位做拒绝采样；再做跨风格
迁移，用\textbf{Bures-Wasserstein 最优传输}同时对齐目标风格的均值与协方差，然后沿
源谱到目标谱的连线做\textbf{单调二分投影}，把缺陷化合物逐样本钳制在物理下界与源
样本值之间。结果表明，条件流的\textbf{负对数似然为 -19.030}，明显优于高斯收缩模型的
10.611；16 种风格的生成样本\textbf{风格保真度均值 0.81}（随机基线 0.045）、多样性比
0.75，未出现模式坍塌；在 Pils/Lager 到 Tripel 的迁移中，被判为 Tripel 的比例由
0.00 升至 \textbf{0.64}，平均迁移强度 $\lambda=0.80$，而\textbf{缺陷 OAV 构造性
不恶化}（DMS 为 1.00 倍、双乙酰 0.84 倍、4VG 0.94 倍），平均缺陷分由 0.113 降至 0.088。

\textbf{创新点}：一是发现并量化了附件未说明的量纲混用（33 列原始浓度）与检出限左
删失（102 列、占 12.60\% 观测），把数据契约固化为三问共享的判定；二是把共表达图
同时用于预测侧的特征增广与归因侧的社区一致性判据，使模块级信号成为可检验的对象；
三是用五路秩融合、双稳定性指标、四判据化学合理性与外部矩阵独立验证，把归因从排行榜
升级为可复核论证；四是把跨风格迁移写成\textbf{带单调缺陷约束的最优传输}，以单一标量
$\lambda$ 刻画迁移强度，兼具可解释性与可审计性。

\vspace{4pt}
\noindent\textbf{关键词：}共表达图；图扩散；多任务学习；稳定性选择；归一化流；
最优传输；啤酒风味
